% !TeX program = xelatex
% UTF-8 transcription of all four source pages, in handwritten page order.
% images/pic01.jpg = 20261012_195124161.jpg
% images/pic02.jpg = 20261012_195129710.jpg
% images/pic03.jpg = 20261012_195134811.jpg
% images/pic04.jpg = 20261012_195139125.jpg
% Source wording, abbreviations, and English spellings are retained.
% Compile standalone with XeLaTeX, or input into a XePersian document.
\ifdefined\settextfont
  \def\SyncopeStandalone{0}
\else
  \def\SyncopeStandalone{1}
  \documentclass[12pt,a4paper]{article}
  \usepackage[margin=2cm]{geometry}
  \usepackage{amsmath,amssymb}
  \usepackage{xepersian}
  \IfFontExistsTF{Amiri}{\settextfont{Amiri}}{\settextfont{FreeSerif}}
  \setlatintextfont{FreeSerif}
\fi
\providecommand{\en}[1]{\lr{#1}}
\newcommand{\Syncopepage}[2]{%
  \clearpage
  \section*{موضوع: براونوالد ـ ۷۱ ـ #2 و سنکوپ}
  {\footnotesize تاریخ: \en{/ /}\quad صفحهٔ #1\par}}
\ifnum\SyncopeStandalone=1 \begin{document}\fi

\Syncopepage{۱}{\en{$\mathrm{BP}\downarrow$}}
\begin{itemize}
\item سنکوپ: ناگهانی، گذرا و کامل؛ \en{LOC} + آتونی؛ بازگشت کامل.
\item افت جریان خون مغز برای ۱۰ ثانیه سنکوپ ایجاد می‌کنه.
\item سن شایع سنکوپ:
  \begin{itemize}
  \item ۱۰ تا ۲۰ سالگی.
  \item ۶۰ تا ۸۰ سالگی.
  \end{itemize}
\item شایع‌ترین علت سنکوپ در هر سن: رفلکس مدیتیت‌ها؛ در همه اول \en{NMS}.
  \begin{itemize}
  \item علت دوم علل قلبی شامل تاکی‌آریتمی‌ها و برادی‌آریتمی‌ها.
  \end{itemize}
\item اورتواستاتیک \en{intolerance}: علائم و نشانه‌ها، اختلال در تنظیم \en{BP}.
\item اورتواستاتیک هایپوتنشن:
  \begin{itemize}
  \item افت \en{20 mm} در \en{SBP} یا \en{10 mm} در \en{DBP}، طی ۳ دقیقه پس از ایستادن.
  \item می‌تواند همراه با علائم باشد یا نباشد.
  \end{itemize}
\item انیشیال اورتواستاتیک هایپوتنشن (فیزیولوژیک است):
  \begin{itemize}
  \item زیر \en{40 mm} است؛ ظرف \en{30 s}، \en{NL} می‌شود.
  \end{itemize}
\end{itemize}

\Syncopepage{۲}{\en{$\mathrm{BP}\downarrow$}}
\begin{itemize}
\item \en{delayed progressive orthostatic Bp hypotention}
  \begin{itemize}
  \item پس از ایستادن به صورت تدریجی \en{$\mathrm{BP}\downarrow$}.
  \end{itemize}
\item داروها شایع‌ترین علت اورتواستاتیک هایپوتنشن هستند.
\item دسته‌بندی سنکوپ‌های رفلکس مدیتیت:
  \begin{itemize}
  \item اگر افت \en{BP} به علت وازودیلیشن محیطی غالب باشه:
    \en{Vasodepressor-type reflex response}.
  \item اگر برادی‌کاردی غالب باشه: \en{Cardioinhibitory response}.
  \item \en{Mixed response}.
  \end{itemize}
\item ازدیاد حساسیت سینوس کاروتید (\en{CSH}):
  \begin{itemize}
  \item علت \en{$1/3$} تمام سنکوپ‌ها در افراد مسن.
  \item با مکانیسم \en{Cardioinhibitory}.
  \item پس از تشخیص اگر سنکوپ تکرار می‌شد: پیس‌میکر دوحفره‌ای.
  \item پس از ماساژ کاروتید:
    \begin{itemize}
    \item آسیستول \en{$>3\,\mathrm{s}$}.
    \item بلاک \en{AV}.
    \item \en{$\mathrm{SBP}\downarrow\;>50\,\mathrm{mm}$}.
    \item تکرار علائم بالینی بیمار.
    \end{itemize}
  \end{itemize}
\end{itemize}

\Syncopepage{۳}{افت \en{BP}}
\begin{itemize}
\item شایع‌ترین تاکی‌آریتمی که سنکوپ می‌دهد: \en{VT}.
\item برخورد با بیمار سنکوپ:
  \begin{itemize}
  \item شرح حال و معاینه: علت ۲۵\% موارد یافت می‌شود.
  \item با شرح حال: آیا واقعاً سنکوپ بوده است؟
  \end{itemize}
\item سنکوپ به دنبال \en{VT} یا \en{AVB} معمولاً پرودروم ندارند و زیر ۵ ثانیه ریکاوری یافته.
\item چک \en{VS} در شرایط اورتواستاتیک:
  \begin{itemize}
  \item پس از ایستادن، در دقیقه تا ۳ دقیقه.
  \item \en{Post Ortho Tachy Syndrom (POTS)}:
    \begin{itemize}
    \item افزایش \en{HR}، ۳۰ تا، بدون افت \en{BP}.
    \end{itemize}
  \end{itemize}
\item پس از چک نبض: ماساژ سینوس کاروتید.
  \begin{itemize}
  \item ۱۰–۱۵ ثانیه در دو طرف، هم ایستاده هم خوابیده.
  \item اگر سابقهٔ \en{TIA} در ۳ ماه اخیر: انجام ندهیم.
  \item پاسخ \en{NL}: افت گذرا \en{HR}، \en{AV prolong}.
  \end{itemize}
\item هیچ آزمایش به صورت روتین چک نمی‌شود.
\item تست تیلت:
  \begin{itemize}
  \item بیمار ۲۰ دقیقه خوابیده؛ با زاویهٔ ۷۰ درجه سپس ۳۰–۴۵ دقیقه ایستاده.
  \item هدف: تأیید تشخیص \en{NMS}.
  \end{itemize}
\end{itemize}

\Syncopepage{۴}{افت \en{BP}}
\begin{itemize}
\item اکو تنها در صورت شک به اختلالات ساختاری قلب:
  \begin{itemize}
  \item اگر \en{EKG}، \en{P/E}، \en{NL} باشه نباید بگیریم.
  \end{itemize}
\item سنکوپ در \en{ACS} ناشایع است، مگر شدید باشه.
\item سنکوپ حین فعالیت: علت قلبی.
\item سنکوپ پس از فعالیت: \en{\fbox{NMS}}.
\item همهٔ سنکوپ‌ها \en{EKG} گرفته شود.
  \begin{itemize}
  \item هرگونه اختلال: بررسی.
  \end{itemize}
\item اگر سنکوپ غیرکامپلیکه است: \en{CT} یا \en{MRI} نمی‌گیریم.
\end{itemize}

% Repeated subject/date labels and handwritten page numbers are reproduced above.
% The printed notebook emblems contain no text.
\ifnum\SyncopeStandalone=1
  \let\SyncopeEnd\enddocument
\else
  \let\SyncopeEnd\relax
\fi
\SyncopeEnd
